\section{Conclusion}\label{sec:conclusion}

This article has shown that a single ABI
can be used for different implementations of a blockchain peer.
This allows one to use the same testcases, unchanged, against different
configurations of the blockchain used for testing, with different consensus
and number of peers, both local and remote. Moreover, a stubbed implementation
of the peer allows one to execute the testcases in a few minutes, which
makes them adequate for continuous integration.

The ability to run the testcases against blockchains with completely
diffent consensus has been invaluable when we implemented
Hotmoka over the Mokamint engine (proof of space), after its
original implementation over the Tendermint engine (proof of stake) had been completed.
In this case, the testcases have been used as non-regression tests.
In general, testcases can both support confidence in the
correctness of the smart contracts and provide a non-regression check for the
peer used for testing.

This article focused on the Hotmoka blockchain and on Takamaka smart contracts,
since it originated during the development of the latter blockchain.
However, a single ABI could be used also for other blockchains, such as Web3 for Ethereum.

A limit of testing for smart contracts is that the blockchain used for
testing must have deterministic finality, since otherwise there might be
history changes that make some tests fail (\cref{fig:experiments}).
This is a general result, not related to the use of a single ABI.
A possible solution is to introduce a delay between transaction requests,
in order to give the blockchain the time to stabilize. But this would
further increase the time for executing the testcases, that is already significant
for non-stubbed peer implementations (\cref{fig:experiments}).
